% Generated by tools/build_cv.py from the career repo (cv/cv.yaml). Do not edit.

\cvsection{Research Projects}

\begin{cventries}

  \cventry
    {PhD Thesis project}
    {Large-scale Mobility-on-demand: Simulation Studies and Optimization}
    {CTU in Prague}
    {2017--2024}
    {
      \begin{cvitems}
        \item {This thesis focuses on improving the large-scale Mobility-on-demand (MOD) systems. Simulation studies are used to better understand MoD systems and answer related research questions. Based on the insights gained from the simulations, we propose a new optimization method for optimal plan chaining. This new method has the potential to improve the efficiency of various MoD-related problems, which we demonstrate with the example of vehicle dispatching (Dial-a-ride problem).}
        \item {Currently, the thesis is submitted and waiting for defense. The last manuscript from the thesis focused on optimal plan chaining is in review in \emph{Transportation Research Part C: Emerging Technologies}. Finally, a subsequent manuscript focused on the plan chaining application is in review in \emph{Nature Computational Science}}
      \end{cvitems}
    }

  \cventry
    {Fulbright project}
    {Designing Networks for Mass Transit Systems Supplemented by Mobility-on-Demand Services}
    {Cornell University}
    {September 2025--June 2026}
    {
      \begin{cvitems}
        \item {Extending mass transit network design so that the evaluation of a network also covers a mobility-on-demand (MoD) service, i.e. the demand not served by the transit lines is served by ridesharing vehicles and their operation cost enters the design objective.}
        \item {Areas with high demand are served by transit lines, low-demand areas and the first/last mile by MoD, so that the mix can be designed automatically from historical travel demand instead of by expert judgement.}
        \item {Joint work with the group of Prof. Samitha Samaranayake at Cornell University, combining our experience with large-scale dial-a-ride optimization with their transit network design expertise.}
      \end{cvitems}
    }

  \cventry
    {Bell Textron Inc. Multi-Modal Traveler Modeling Project}
    {\hrefs{https://www.aic.fel.cvut.cz/news/air-taxis-as-the-future-of-public-transport}{Feasibility of eVTOL Transport in Urban Environments}}
    {CTU in Prague}
    {2021}
    {
      \begin{cvitems}
        \item {This project studied the viability of using electric vertical take-off and landing travel system for urban mobility.}
        \item {Our transportation team defined the goals and requirements with Bell Textron Inc. and developed a simulation model to evaluate the feasibility of the eVTOL system.}
        \item {I focused on the simulation model and optimal eVTOL station positioning.}
        \item {We delivered the codebase and summarized our findings in a report.}
      \end{cvitems}
    }

\end{cventries}
